% Generated draft: requires mathematical and source review.
\documentclass[11pt]{article}
\usepackage[margin=25mm]{geometry}
\usepackage{fontspec}
\setmainfont{Latin Modern Roman}
\usepackage{amsmath,amssymb,mathtools,unicode-math}
\setmathfont{Latin Modern Math}
\usepackage{enumitem,needspace,xurl,hyperref}
\setlength{\parindent}{0pt}
\setlength{\parskip}{5pt}
\setlength{\emergencystretch}{3em}
\newcommand{\N}{\mathbb{N}}
\newcommand{\Z}{\mathbb{Z}}
\newcommand{\Q}{\mathbb{Q}}
\newcommand{\R}{\mathbb{R}}
\newcommand{\C}{\mathbb{C}}
\newcommand{\F}{\mathbb{F}}
\newcommand{\A}{\mathbb{A}}
\newcommand{\PP}{\mathbb{P}}
\newcommand{\E}{\mathbb{E}}
\newcommand{\eps}{\varepsilon}
\newcommand{\dd}{\,\mathrm{d}}
\begin{document}
\section*{UnsolvedMath ID: 700707}
\textbf{Source problem code:} catalogue entry
\subsection{Definitions and mathematical statement}
Let \(k\) be an algebraically closed field of characteristic \(2\). For a finite group \(G\), let \(S_1,\ldots,S_r\) represent the isomorphism classes of simple \(kG\)-modules, and let \(P_j\) be the projective cover of \(S_j\). Define the Cartan matrix of the full group algebra by
\[
C_k(G)=(c_{ij})_{1\leq i,j\leq r},\qquad c_{ij}=[P_j:S_i],
\]
where \([P_j:S_i]\) is the multiplicity of \(S_i\) among the composition factors of \(P_j\). Write \(C_k(G)\sim C_k(H)\) if the matrices agree after a simultaneous permutation of their rows and columns, corresponding to a bijection between the simple-module isomorphism classes.

Under the interpretation that characterization means determination among all finite groups, the question is whether, for every algebraically closed field \(k\) of characteristic \(2\),
\[
\forall S\text{ finite simple}\;\forall G\text{ finite},\qquad
C_k(G)\sim C_k(S)\ \Longrightarrow\ G\cong S.
\]
Here a simple group is nontrivial and has no nontrivial proper normal subgroup. The supplied wording does not explicitly state whether comparison is among all finite groups or only among finite simple groups. Under the latter interpretation, the quantifier over \(G\) is restricted to finite simple groups. The full group-algebra Cartan matrix and the simultaneous-permutation convention above are standard interpretations that require confirmation against the original source.

As background, the supplied record attributes to R. Brandl's 1982 paper the result that, for any given finite group, only finitely many finite groups, counted up to isomorphism, have the same Cartan matrix. This is a reported source claim; its precise hypotheses and matrix-equivalence convention have not been independently checked.
\subsection{Short English statement}
Does the Cartan matrix in characteristic \(2\) determine every finite simple group up to isomorphism? The supplied wording leaves unclear whether the comparison groups may be arbitrary finite groups or must also be simple.
\subsection{Sources}
\begin{itemize}
\item Local UnsolvedMath registry, permanent ID \(700707\), exact source problem code \texttt{catalogue entry}, title \emph{Kourovka Notebook Problem 11.14}. The supplied registry labels the problem open; this label is not an independently verified assessment.
\item \emph{The Kourovka Notebook (Unsolved Problems in Group Theory)}, Problem 11.14, proposed by R. Brandl in the 11th issue (1990). The supplied background identifies the 21st issue edition (2026), \texttt{arXiv:1401.0300v48}. The registry describes the collection as edited by E. I. Khukhro and V. D. Mazurov.
\item R. Brandl, \emph{Arch. Math.} \textbf{38} (1982), 322--323. The supplied statement cites this paper for the finiteness assertion recorded above. No article title or further bibliographic details are supplied.
\end{itemize}
\end{document}
